\documentclass{article}
\usepackage{graphicx}
\usepackage[margin=1.5cm]{geometry}
\usepackage{amsmath}

\begin{document}
\twocolumn

\title{Wednesday Warm Up: Unit 5, Electromagnetism and Optics}
\author{Prof. Jordan C. Hanson}

\maketitle

\section{Memory Bank}

\begin{itemize}
\item $c/n = f \lambda$ ... Relationship between the speed of light, $c$, the index of refraction, $n$, the frequency of an electromagnetic wave, $f$, and its wavelength, $\lambda$.
\item $c = 3.0 \times 10^8$ m s$^{-1}$ ... The speed of light in a vacuum, within 1 percent error.
\item $\Delta x = v \Delta t$ ... The relationship between displacement, constant velocity, and time duration.
\item \textbf{Optics:} Snell's Law states that, for an interface between two media with \textit{indices of refraction} $n_1$ and $n_2$, light leaves one medium with an angle $\theta_1$ and enters the second medium with an angle $\theta_2$ as follows:
\begin{equation}
n_1 \sin\theta_1 = n_2\sin\theta_2
\end{equation}
\noindent The two angles are defined with respect to vertical.
\item \textbf{Optics:} Fresnel's equations give the fraction of power (in light) that is reflected ($R$) or transmitted ($T$) by an interface between media with indices of refraction $n_1$ and $n_2$:
\begin{equation}
R = \left| \frac{n_1 - n_2}{n_1 + n_2} \right|^2 \label{eq:2}
\end{equation}
\noindent Note that, to conserve power, $T = 1 - R$.
\end{itemize}

\section{Electromagnetic Waves}

\begin{enumerate}
\item The frequency of a stream of monochromatic optical photons is $4.3 \times 10^{14}$ Hz.  (a) What is the wavelength of this light in air ($n \approx 1$)?  (b) What is the wavelength in glass, with $n = 1.5$? \\ \vspace{2cm}
\item Suppose we send an RF signal down a coaxial cable that is 900 meters long.  (a) If we observe the signal at the other end of the cable after 4 $\mu$s, what is the speed of the RF pulse in the cable? (b) What is the index of refraction of the cable? \\ \vspace{2cm}
\item A 60-kW radio transmitter on Earth sends its signal to a satellite 100 km away.  At what distance in the same direction would the signal have the same maximum field strength if the transmitter’s output power were increased to 90 kW? \\ \vspace{2cm}
\end{enumerate}

\section{Light Can Refract and Reflect}

\begin{enumerate}
\item Suppose a ray of sunlight hits a portion of the ocean that is momentarily still and flat.  If the sun is 45 degrees away from vertical, at what angle (from vertical) do the rays enter the water?  Assume the index of refraction of air is $n_1 = 1.0$, and that of the water is $n_2 = 1.336$. \\ \vspace{3cm}
\item Suppose a ray of sunlight hits a portion of the ocean that is momentarily still and flat.  (a) If the sun is 0 degrees away from vertical, what fraction of power is reflected?  Assume the index of refraction of air is $n_1 = 1.0$, and that of the water is $n_2 = 1.336$. (b) What fraction of power is transmitted into the ocean? \\ \vspace{3cm}
\item Think about the last time you were at the beach.  Does the image of the Sun on the water look brighter when the sun is \textit{lower} or \textit{higher} in the sky?  How should Eq. \ref{eq:2} be modified (qualitatively) to account for this? \\ \vspace{2cm}
\end{enumerate}

\end{document}
